% ================================================================
%  chapters/parts/ch03-a.tex
%  第 3 章  热分析方法（分片 a）
%  源：S1（上册）第 3 章，书内 p.169—180 / PDF p.188—199
%  覆盖：§3.1 热分析的分类、§3.2 几种常用的热分析技术（TG / DTA / DSC / MTDSC）、
%        §3.3 热分析动力学简介（基本关系与等转化率法的前提）
%  未覆盖（由后续分片续写）：§3.4 应用、§3.5 联用技术、§3.6 实验技巧、§3.7 结束语
% ================================================================

\chapterstart{热分析方法}%
{本章回答：热分析曲线上的台阶与峰究竟对应哪些物理量，热重法（TG）、差热分析（DTA）与
差示扫描量热法（DSC）的信号含义与纵轴正负约定为何不同，什么条件下才能把一条曲线换算成
组成与热量，以及热分析求出的活化能在什么前提下才可信。§3.1 给出分类与三种主力技术的分工，
§3.2 逐个技术讲原理与读图，§3.3—§3.6 沿「测量与条件—数据到结论—误判与边界—互补与回查」
这条主线展开，§3.7 讲联用热分析（源书 §3.4—§3.7 的应用、联用、实验技巧与结束语四节已按这条主线重编）。}
\label{ch:03}

\subsection{这个方法测量什么}

\begin{probout}
热分析（thermal analysis，TA）的定义里有两个限定条件必须与「测量」二字一起读：
\textbf{程序控温}与\textbf{一定气氛}。它测量的是物质的某一物理性质随温度或时间的变化，
输出只有一条曲线，曲线本身不给出组成、结构或机理，只回答「在什么温度发生了什么变化」
\src{S1 §3.1，书内 p.171 / PDF p.190}。

按所测的物理性质，ICTA 把热分析分为九类十七种：质量（热重法）、温度（差热分析）、热量
（差示扫描量热法、温度调制式差示扫描量热法）、尺寸与力学量（热机械分析、动态热机械分析）、
声学量、光学量、电学量、磁学量，另有把两种以上技术接在一起的联用技术——TG-DTA、TG-DSC 属
同时联用，TG/MS、TG/FTIR 属串接联用，TG/GC 属间歇联用
\src{S1 §3.1 表 3-1，书内 p.171--172 / PDF p.190--191}。催化剂研究中真正用得多的只有三种，
它们的区别在\textbf{测的物理量本身}：TG 测质量、DTA 测温差、DSC 测热流（能量差）。
\end{probout}

「TG 测质量、DTA 测温差、DSC 测热流」是本章最需要记住的一句话：三者的曲线形状可以很相似，
信息含量却相差很大，而混淆往往从纵坐标开始。表~\ref{tab:ch03-three} 按五个问题把三者拆开：
测什么、纵轴方向怎么定、能给出什么、不能给出什么、定量能力到哪一步。

\begin{table}[tbp]
\centering
\caption{热重法、差热分析与差示扫描量热法的对照。读者应看两行：第二行的纵轴方向约定
（DTA 与 DSC 相反，是读图时最容易弄反的地方），以及最后两行的能力边界——差异的根源在于
DTA 的信号是温差、DSC 的信号是热流，二者之间隔着一个需要标定的热阻。\src{据 S1（上册）
§3.2.2—§3.2.3，书内 p.171--177 / PDF p.190--196 归纳；原书表 3-1 为「热分析技术分类」
（ICTA 将热分析技术分为 9 类 17 种），与本表无对应关系，故此处只引正文。}}
\label{tab:ch03-three}
{\srcsize
\begin{tabularx}{\textwidth}{@{}lXXX@{}}
\toprule
项目 & 热重法 TG & 差热分析 DTA & 差示扫描量热法 DSC \\
\midrule
直接测量的物理量 & 试样质量 $m$，信号为变化量 $\Delta m$ & 试样与参比物的温度差 $\Delta T$ & 输给试样与参比物的能量差，即热流率 $\mathrm{d}Q/\mathrm{d}t$ \\
纵轴方向约定 & 向上为质量增加、向下为质量减小；DTG 失重峰朝下 & 习惯约定：向上放热、向下吸热 & 热力学惯例：向上为正＝吸热、向下为负＝放热 \\
能给出什么 & 台阶的温度区间与失重百分数、残余量、起始分解温度（表征热稳定性）；由台阶高度反推组成 & 热事件的温度与先后次序、吸热或放热的性质；ICTA 规定用外推起始温度表示反应起始温度 & 热事件的定量热焓 $\Delta H$（峰面积）、比热容、转化率—温度（时间）关系；经标定后可溯源 \\
不能给出什么 & 无质量变化的相变（熔化、晶型转变）；逸出气体的种类；同一温区内重叠失重过程的分辨 & 热量：峰面积的比例系数含装置热阻与样品热容，未经标定不能积分成 $\Delta H$；峰温不能作鉴定物质的特征温度 & 无热效应的质量变化；逸出气体的种类；定量准确度依赖温度与热量两项标定 \\
定量能力 & 定量（质量可直接溯源到天平） & 通常只能定性或半定量 & 定量，峰面积可换算为 $\unit{J.g^{-1}}$ \\
\bottomrule
\end{tabularx}}
\end{table}

三者的差别不只是「测的量不同」，而决定了各自能支撑什么结论。TG 的纵坐标是质量，质量可以
直接溯源到天平，台阶高度就是失重百分数；DSC 的纵坐标是热流率，经温度与热量两项标定后峰
面积可以换算成热焓，也是定量技术。DTA 的纵坐标是温差 $\Delta T$，而 $\Delta T$ 与热流率
之间隔着一个热阻：同样的放热速率，在不同装置、不同样品量、不同温度下产生的 $\Delta T$
可以相差数倍。未经逐条件标定的 DTA 峰面积不能直接换算成热量——这正是 DSC 被发展出来的
原因，也是「DSC 能做定量、DTA 通常只能定性」的根本差别
\src{S1 §3.2.3，书内 p.177 / PDF p.196}。读图前必须先看纵轴标注与方向约定，否则会把分解
（吸热）与燃烧（放热）的归属整个弄反。

\subsection{核心原理与公式}

\begin{principle}
三种技术的信号都来自试样与环境的能量交换，但取信号的环节不同。TG 取的是\textbf{质量}这一
状态量，信号不经过任何热传递假设，因此最直接；DTA 取的是试样与参比物的\textbf{温差}，
DSC 取的是为维持二者温度相等而输入试样的\textbf{补偿功率}（热流率）。后两者都要经过热传递
环节：DSC 用补偿把这一环节的系数固定下来并加以标定，DTA 没有这一步，所以同一个热事件在
两种仪器上的峰高不可比、在 DTA 上甚至不可积分。
\end{principle}

\subsubsection{热重法（TG）与微商热重（DTG）}

TG 的仪器通称热天平，基本单元是微量电子天平（或石英微天平）、炉子、温度程序器、气氛控制
器与记录系统。记录方式值得注意：仪器先同时存储质量、温度与时间三组数据，测量结束后再由
时间换算成温度 \src{S1 §3.2.1，书内 p.174 / PDF p.193}。这意味着温度滞后会直接进入数据，
而不只是影响显示。

TG 曲线的读法只有两条规则：\textbf{水平段（平台）表示质量恒定，斜率发生变化的段落表示质量
变化}；纵坐标向上表示质量增加、向下表示质量减小。微商热重曲线（DTG 曲线）是质量变化率
$\mathrm{d}m/\mathrm{d}t$ 对温度或时间的图，质量减小时峰朝下。DTG 与 TG 的对应关系是严格的：
DTG 峰的起止点就是 TG 台阶的起止点，峰的数目等于台阶的数目，峰位（最大失重速率处）与 TG
曲线的拐点相对应 \src{S1 §3.1，书内 p.173 / PDF p.192；§3.2.1，书内 p.175 / PDF p.194}。
曲线上可直接读出的量有四个：台阶高度 $\Delta m$、质量变化速率 $\mathrm{d}m/\mathrm{d}t$、
质量开始变化的起始分解温度 $T_{\mathrm{i}}$（用来表征样品的热稳定性）与 DTG 峰温
$T_{\mathrm{p}}$。\textbf{这里的「四个」按原页面图像核对}（书内 p.175 / PDF p.194）：原书在同一段
逐条定义 $\Delta m$、$\mathrm{d}m/\mathrm{d}t$、$T_{\mathrm{i}}$、$T_{\mathrm{p}}$ 四项；
原书记质量用 $W$（该页前文印 $\mathrm{d}W/\mathrm{d}t$ 与 $\mathrm{d}^{2}W/\mathrm{d}t^{2}=0$），
同段的定义句又写 $m$，译文取 $m$。

用失重分数（转化率）$\alpha$ 可以把台阶高度写成与取样量无关的形式：

\begin{equation}\label{eq:ch03-alpha}
\alpha=\frac{m_0-m_t}{m_0-m_\infty}
\end{equation}

式中 $m_0$ 为初始质量，$m_t$ 为温度 $T$（或时间 $t$）时的质量，$m_\infty$ 为过程结束后的
残余质量。$\alpha=0$ 对应未反应，$\alpha=1$ 对应过程结束。

\textbf{这个式子的来源。}此处已对照原页面图像（书内 p.174--175 / PDF p.193--194）：原书 §3.1
只逐条定义 $\Delta m$（质量的变化量）、$\mathrm{d}m/\mathrm{d}t$（质量变化的速率）、
$T_{\mathrm{i}}$（起始分解温度）与 $T_{\mathrm{p}}$（DTG 峰温），\textbf{没有}给出以 $m_0$、
$m_t$、$m_\infty$ 表示的转化率定义式；原书直到 §3.3 的动力学方程才用到 $\alpha$，并只在变量
说明里写「$\mathrm{d}\alpha$ 为转化率」（书内 p.180 / PDF p.199）。因此本式是改编本补出的
定义（量的含义与通用约定一致），用它可以不依赖取样量地表示台阶高度。

\begin{keycond}
式~\eqref{eq:ch03-alpha} 要求 $m_\infty$ 有确定值，即曲线上存在明确的终止平台。若后续温区
仍有失重，或相邻台阶彼此重叠，$m_\infty$ 就成了人为选定的数，$\alpha$ 随之失去意义。此时
应改用绝对失重百分数 $\Delta m/m_0$，并写明用于确定基线（平台）的温区。TG 上真正能定量的
是台阶高度，能不能把它换算成组成，取决于终止平台是否真实存在，而不取决于天平的分辨率。
\end{keycond}

\subsubsection{差热分析（DTA）}

DTA 的原理是「同炉同时加热、比较两个温度」：试样与参比物处于相同的加热或冷却条件下，两支
测温热电偶分别插在试样与参比物中，输出二者随温度变化的温差 $\Delta T$。试样无变化时二者
温度相等、$\Delta T=0$；试样发生吸热或放热变化时二者温度不等，输出温差信号，经放大后记录，
这就构成了 DTA 的基本原理 \src{S1 §3.2.2，书内 p.175 / PDF p.194}。参比物必须在所测温区内
不发生热效应，这是 DTA 实验成立的前提，也是测量与条件一节要单独交代的原因。

DTA 曲线上的特征量在原书中逐条给出：\textbf{基线}（曲线上 $\Delta T$ 近似等于 $0$ 的区段，
如 $abd$）、\textbf{峰}（曲线离开基线又回到基线的部分，峰向下为吸热峰，向上则相反）、
\textbf{峰宽}（曲线偏离基线又返回基线两点间的距离或温度间距，如 $ab$）、\textbf{峰高}
（试样和参比物之间的最大温度差，即峰顶至内插基线间的垂直距离，如 $ce$）、
\textbf{峰面积}（峰和内插基线之间所包围的面积，如 $abc$）、\textbf{外推始点}
（峰的起始边陡峭部分的切线与外延基线的交点，如 $T_{\mathrm{ei}}$ 点）、
\textbf{峰温} $T_{\mathrm p}$
\src{S1 §3.2.2，书内 p.176--177 / PDF p.195--196}。其中两条结论直接决定读图习惯。

\textbf{峰温不能当特征温度用。}峰顶温度只表示试样与参比物温差最大的一点，它既不等于反应的
终止温度，位置又受试样条件影响较大，因此一般不能作为鉴定物质的特征温度，只有在试样条件
相同时才可作相对比较 \src{S1 §3.2.2，书内 p.177 / PDF p.196}。

\textbf{起始温度用外推始点。}ICTA 对大量试样的测定结果表明，外推起始温度与其它实验测得的
反应起始温度最为接近，因此规定用外延起始温度表示反应的起始温度（同一页）。这与上一条互为
补充：一个位置（峰温）不可作特征温度，另一个位置（外推始点）才是可比的起点。

\textbf{峰面积与热量之间差一个比例系数}：

\begin{equation}\label{eq:ch03-dta}
\Delta H=\frac{K}{m}\int_{t_1}^{t_2}\Delta T\,\mathrm{d}t
\end{equation}

式中 $K$ 为比例系数，与装置热阻、样品热导率、热容及温度有关；$m$ 为试样质量。$K$ 不由几何
量决定，只能逐条件标定。未标定时式~\eqref{eq:ch03-dta} 只能用于同一装置、同一条件下不同
样品之间的相对比较。

此处已对照原页面图像（书内 p.176 / PDF p.195 与书内 p.177 / PDF p.196）：原书 §3.2.2 只逐条定义
了 DTA 曲线的基线、峰、峰宽、峰高、峰面积、外推始点与峰温（上一段已按原页补齐），
\textbf{并未给出峰面积与热量之间的换算关系}，原书全节没有出现符号 $K$、$m$，也没有积分式；
式~\eqref{eq:ch03-dta} 是本改编本补出的通用定量关系（峰面积正比于热效应，比例系数只能逐条件
标定），不是原书内容。原书此处的配图为图~3-3（热电偶和温差热电偶）与图~3-4（典型 DTA
曲线）。另：原书该段自身在「外推起始温度」与「外延起始温度」两种写法间混用，译文照录。

\subsubsection{差示扫描量热法（DSC）}

DSC 是为克服 DTA 在定量测定上的不足而发展起来的。它测量与试样热容成比例的单位时间功率输出
与程序温度或时间的关系，做法上完成了 DTA 做不到的一件事：\textbf{及时补偿试样因热效应产生
的能量变化，使试样与参比物之间的温度始终保持相同}——无温差、无热传递，热损失小、检测信号
大，因此灵敏度与精度都大幅提高，可以进行热量的定量分析
\src{S1 §3.2.3，书内 p.177 / PDF p.196}。

DSC 量热仪分两类。\textbf{功率补偿型}采用零点平衡原理，又分外加热与内加热两种：外加热式把
试样和参比物放在外加热炉内加热，二者容器下各装一组独立的补偿加热丝与传感器，两个控制系统
分别负责按预定速率升降温、以及补偿二者之间产生的温差（试样出现温差时改变补偿加热丝的电流）；
内加热式不设外加热炉，直接用两个小加热器加热并同时进行功率补偿。\textbf{热流型}仍按差热
分析的原理做量热，分热流式与热通量式；热通量式的检测器由许多热电偶串联成热电堆式热流量计，
两个热流量计反向连接，分别安装在试样容器、参比容器与炉体加热块之间，热端均匀分布在容器壁
上，因此检测信号大、测得的试样温度是试样各点温度的平均值，重复性、灵敏度与精确度都高，
常用于精密的热量测定 \src{S1 §3.2.3，书内 p.177--178 / PDF p.196--197}。

DSC 曲线的纵坐标是热流率 $\mathrm{d}Q/\mathrm{d}t$，单位 $\unit{mW}$（即
$\unit{mJ.s^{-1}}$），横坐标是温度或时间；按热力学惯例，向上为正、表示吸热效应，向下为负、
表示放热效应 \src{S1 §3.1，书内 p.173 / PDF p.192}。峰面积给出热焓：

\begin{equation}\label{eq:ch03-dsc}
\Delta H=\frac{1}{m}\int_{t_1}^{t_2}
\left[\frac{\mathrm{d}Q}{\mathrm{d}t}
-\left(\frac{\mathrm{d}Q}{\mathrm{d}t}\right)_{\mathrm{base}}\right]\mathrm{d}t
\end{equation}

积分区间从峰的外推始点到外推终点，括号内是被测信号与基线之差。与式~\eqref{eq:ch03-dta}
相比，关键差别在于热流率与功率之间的关系由仪器标定给出，不必像 DTA 那样每次实验重新确定
一个与装置和样品都相关的 $K$，所以峰面积可以直接换算成 $\unit{J.g^{-1}}$。

\textbf{原书 DSC 曲线的「数学表达式」。}此处已对照原页面图像（书内 p.178 / PDF p.197）：原书
§3.2.3 的「2. 差示扫描量热曲线」只印一句 $\mathrm{d}H/\mathrm{d}t=f(T)$ 或 $f(t)$，并说明以
热流率为纵坐标的曲线与 DTA 相似、不再重复，\textbf{没有}给出峰面积与热焓之间的换算式；原书把
热流率记作 $\mathrm{d}H/\mathrm{d}t$，本稿统一记作 $\mathrm{d}Q/\mathrm{d}t$（纵坐标
$\unit{mW}$ 的口径不变）。式~\eqref{eq:ch03-dsc} 是改编本补出的定量关系。

\begin{keycond}
式~\eqref{eq:ch03-dsc} 成立需要三个条件：一是温度与热量两项标定在测量前完成，且在测量温区
内有效；二是基线由同条件下的空白（空坩埚、同升温速率、同气氛）确定，不能凭肉眼在两段平台
之间任意连线；三是峰的外推始点与外推终点可辨识，重叠峰必须先分峰，否则应改用其它方法。
灵敏度提高不等于准确度提高：微量样品热容小、信噪比好，但称量误差与热接触差异会同步放大，
报 $\Delta H$ 时应给出重复测量的次数与分散度。
\end{keycond}

\subsubsection{温度调制式差示扫描量热法（MTDSC）}

MTDSC 在传统线性变温程序上叠加一个正弦振荡温度程序，使热流量与温度信号形成非线性调制，
从而把总热流信号分解成可逆与不可逆两种成分 \src{S1 §3.2.4，书内 p.178 / PDF p.197}：

\begin{equation}\label{eq:ch03-mtdsc}
T(t)=T_0+\beta t+A_T\sin(\omega t),\qquad
\frac{\mathrm{d}Q}{\mathrm{d}t}=C_p\frac{\mathrm{d}T}{\mathrm{d}t}+f(T,t)
\end{equation}

式中 $T_0$ 为起始温度，$\beta$ 为基础（线性）升温速率，$A_T$ 为调制幅度，$\omega$ 为调制
角频率，$C_p$ 为试样热容。第一项与升温速率成正比，经傅里叶变换后落在调制频率上，对应与
热容有关的\textbf{可逆}热流成分（玻璃化转变、大多数熔融）；第二项只与温度和时间有关，对应
动力学控制的\textbf{不可逆}成分（蒸发、挥发、结晶、热固化、分解与部分熔融）。分解效果主要
取决于三个量的选择：基础升温速率、调制周期与调制幅度
\src{S1 §3.2.4，书内 p.179 / PDF p.198}。

\textbf{原书的对应式。}此处已对照原页面图像（书内 p.179 / PDF p.198）：原书印
$\mathrm{d}H/\mathrm{d}t=C_p\,\mathrm{d}T/\mathrm{d}t+f(T,t)$，并说明第一项为可逆部分
（包括热容、玻璃化转变及大多数熔融）、$f(T,t)$ 为不可逆部分（包括焓恢复、挥发、结晶、
热固化、蛋白质变性、淀粉凝胶、分解和一些熔融）；本稿把热流率记作 $\mathrm{d}Q/\mathrm{d}t$，
故第二式与原书逐项同形而符号不同。前半的正弦温度程序
$T(t)=T_0+\beta t+A_T\sin(\omega t)$ 是改编本补出的写法——原书只用文字说明「在传统线性变温
基础上叠加一个正弦振荡温度程序」并给出图 3-8，没有印出该式。

MTDSC 的曲线形状与普通 DSC 完全一样，总热流等同于普通 DSC 的结果，额外得到的是分离后的
两条分量曲线：基础升温速率缓慢使分辨率提高，瞬时升温速率快使灵敏度提高，一次实验同时得到
比热容与热流，这是它的独特之处 \src{S1 §3.2.4，书内 p.179--180 / PDF p.198--199}。需要
注意的是，可逆与不可逆的划分是按\textbf{对升温速率的响应方式}做的操作划分，不是热力学上的
可逆性；同一变化在不同调制参数下可能落入不同分量，因此报告分量曲线必须同时给出这三个参数，
并与普通 DSC 的总热流曲线对照检查。

\subsubsection{热分析动力学的基本关系}

程序控温可分为等温与非等温两种，由于后者具有独特优点，已逐渐成为热分析动力学的核心。一般
所采用的不定温、非均相反应的动力学方程为（原书给出的是以温度为自变量的形式）
\src{S1 §3.3，书内 p.180 / PDF p.199}：

\begin{equation}\label{eq:ch03-rate}
\frac{\mathrm{d}\alpha}{\mathrm{d}t}=k(T)f(\alpha)
\qquad\Longleftrightarrow\qquad
\frac{\mathrm{d}\alpha}{\mathrm{d}T}=\frac{1}{\beta}k(T)f(\alpha)
\end{equation}

式中 $\alpha$ 为转化率（式~\eqref{eq:ch03-alpha}），$\beta=\mathrm{d}T/\mathrm{d}t$ 为升温
速率，$k(T)$ 为速率常数的温度关系式，$f(\alpha)$ 为反应机理函数。速率常数的温度关系取
Arrhenius 形式：

\begin{equation}\label{eq:ch03-arrh}
k(T)=A\exp\!\left(-\frac{E_{\mathrm{a}}}{RT}\right)
\end{equation}

$A$ 为指前因子，$E_{\mathrm{a}}$ 为活化能，$R$ 为气体常量，$T$ 为热力学温度。把
式~\eqref{eq:ch03-arrh} 代入式~\eqref{eq:ch03-rate} 并分离变量积分，得到等转化率法的出发点：

\begin{equation}\label{eq:ch03-integral}
g(\alpha)=\int_0^{\alpha}\frac{\mathrm{d}\alpha}{f(\alpha)}
=\frac{A}{\beta}\int_0^{T}\exp\!\left(-\frac{E_{\mathrm{a}}}{RT}\right)\mathrm{d}T
\end{equation}

\textbf{这两式与原书的关系。}此处已对照原页面图像（书内 p.180--181 / PDF p.199--200）：原书
§3.3 只印以温度为自变量的动力学方程
$\mathrm{d}\alpha/\mathrm{d}T=(1/\beta)k(T)f(\alpha)$（即上式右半，逐字同形），随后给出
Ozawa--Flynn--Wall（$\mathrm{d}(\lg\beta)/\mathrm{d}(1/T)=-0.4567E/R$）、Kissinger
（$\mathrm{d}[\ln(\beta/T_{\mathrm{p}}^{2})]/\mathrm{d}(1/T_{\mathrm{p}})=-E/R$）与
Freeman--Carroll 三种处理方法，\textbf{没有}印 Arrhenius 形式的 $k(T)$，也没有印积分式
$g(\alpha)$；式~\eqref{eq:ch03-arrh} 与式~\eqref{eq:ch03-integral} 是改编本按学科标准形式补出
的两步（原书以 $E$、$k(T)$ 直接进入作图法，未展开这两式）。另：原书该处的变量说明把升温速率
写成「$\mathrm{d}T$」而式中用的是 $\beta$，原书写法记于勘误乙类 B13。

\begin{keycond}
\textbf{（1）单一升温速率下求出的 $E_{\mathrm{a}}$ 不可靠。}式~\eqref{eq:ch03-rate} 中有
$A$、$E_{\mathrm{a}}$ 与机理函数 $f(\alpha)$ 三组未知量，在一条曲线上它们的作用是耦合的：
把 $E_{\mathrm{a}}$ 调大、同时换一个 $f(\alpha)$，可以在同等程度上拟合同一条曲线。因此
单条 TG 曲线给出的「活化能」既依赖所设机理函数，也依赖拟合所选温区，数值可以相差很多。
原书的做法是用 $4$ 个以上不同的升温速率测定同一系列曲线，得到一组随升温速率提高而向高温
推移的 TG 曲线，再由斜率求 $E_{\mathrm{a}}$ \src{S1 §3.3，书内 p.180--181 / PDF p.199--200}。
\textbf{（2）动力学三要素必须同时确定。}$E_{\mathrm{a}}$、$A$ 与 $f(\alpha)$ 共同决定
$\mathrm{d}\alpha/\mathrm{d}t$；只报一个活化能而不给机理函数与指前因子，别人无法复现任何
一条曲线，也无法外推到别的温度程序。\textbf{（3）等转化率法的前提是同一 $\alpha$ 处机理
不变。}若求得的 $E_{\mathrm{a}}$ 随 $\alpha$ 明显变化，说明过程中机理发生了改变，或存在
扩散控制与重叠反应，此时把它当作单一常数使用没有意义。上述三条是方法本身的性质，不是实验
误差；原书 §3.3 未逐条写明，此处按方程组的数学结构补入
\supp{原书 p.180--181 只给出动力学方程与作图法，未展开上述前提与失效判据；本条为编写组按
学科通行判定整理（2026-02），未引外部文献}。
\end{keycond}

\textbf{等转化率法的基本思想。}式~\eqref{eq:ch03-integral} 的左边 $g(\alpha)$ 只与转化率
有关，右边只与温度程序有关。在若干个升温速率 $\beta_1,\beta_2,\dots,\beta_n$ 下各测一条
曲线，取同一转化率 $\alpha$ 处的温度 $T_{\alpha,i}$，则 $g(\alpha)$ 对所有这些实验是同一个
值——于是可以\textbf{在不预设机理函数 $f(\alpha)$ 的前提下把它消去}，由 $\ln\beta_i$ 对
$1/T_{\alpha,i}$ 的线性关系直接求出 $E_{\mathrm{a}}$。Ozawa--Flynn--Wall 法是在固定的
$\alpha$ 处使用这一关系，通常由 $4$ 个以上不同升温速率的 TG 曲线作图；Kissinger 法则取每条
曲线的峰值温度 $T_{\mathrm{p}}$ 作为某个特定转化率所对应的温度，用同类的线性化关系求
$E_{\mathrm{a}}$ \src{S1 §3.3，书内 p.180--181 / PDF p.199--200}。两种方法的具体线性化公式，
以及 Freeman--Carroll 等由单条曲线同时求 $E_{\mathrm{a}}$ 与反应级数的方法，见后文 §3.3 续。

\subsection{测量与条件}

\begin{measure}
\textbf{样品。}样品量决定传热滞后：量大则试样内部温度梯度大、峰变宽、峰温向高温移动、相邻
峰合并。常规筛选用几毫克至十几毫克；动力学测量与重叠峰的分离要求更少，并尽量铺成薄层、
颗粒细小、与坩埚底良好接触。\textbf{坩埚。}常用铂、铝、石英与刚玉坩埚
\src{S1 §3.2.1，书内 p.174 / PDF p.193}；坩埚材料不得与样品或气氛发生反应（见后文误判
第（5）条）。加盖可抑制挥发、提高峰的分辨，但会改变产气过程的自生分压，凡涉及产气的定量
必须写明是否加盖。\textbf{气氛。}TG 可在静态或流动态下进行；静态条件下反应生成的气体会在
试样周围累积，试样反应速率随产物的分压而变，因此一般建议在动态气流下测量。常用气体有
\ce{Ar}、\ce{N2}、\ce{O2}、空气、\ce{CO2}、\ce{H2}、\ce{Cl2}
\src{S1 §3.2.1，书内 p.174 / PDF p.193}。气氛是实验条件的一部分：同一催化剂在空气中测到的
放热峰，在氮气中常常消失。\textbf{升温速率。}$\qtyrange{1}{20}{K.min^{-1}}$ 是常规范围，
速率越快灵敏度越高、分辨率越低、峰温越向高温移动；动力学测量至少需要 $4$ 个升温速率。
\textbf{校准。}温度用高纯金属（\ce{In}、\ce{Sn}、\ce{Zn} 等）的熔点标定，热量用已知熔化焓
的标样或蓝宝石比热标定，TG 用标准砝码标定；DSC 的基线用空坩埚在同程序、同气氛下测空白
确定。\textbf{报告。}升温速率、气氛与流量、样品量、坩埚材料与是否加盖是曲线可比性的前提，
缺任何一项，别人都无法重复这条曲线。
\end{measure}

\subsection{数据到结论}

\begin{figure}[tbp]
\centering
\begin{tikzpicture}
% declare function 要用 \tikzset 在 tikzpicture 层声明，图内所有 axis 才能共用。
% （放进单个 \nextgroupplot 只对该子图有效 → Unknown function；
%   放进 \begin{groupplot}[...] 又不是合法的 group 键 → I do not know the key。
%   两种写法本项目都实际踩过。）
\tikzset{declare function={
  stp(\u,\w)=0.5*(1+\u/sqrt(\u*\u+\w*\w));
  dtp(\u,\w)=0.5*\w*\w/((\u*\u+\w*\w)*sqrt(\u*\u+\w*\w));
  tgm(\x)=100-8*stp(\x-110,12)-14*stp(\x-255,14)-22*stp(\x-355,20)-3*stp(\x-500,14);
  dtm(\x)=-(8*dtp(\x-110,12)+14*dtp(\x-255,14)+22*dtp(\x-355,20)+3*dtp(\x-500,14));
  dsm(\x)=0.02+0.00035*(\x-40)
    +0.55/(1+((\x-112)*(\x-112)/121)*((\x-112)*(\x-112)/121))
    +0.75/(1+((\x-258)*(\x-258)/100)*((\x-258)*(\x-258)/100))
    +0.35/(1+((\x-318)*(\x-318)/784)*((\x-318)*(\x-318)/784))
    -0.60/(1+((\x-372)*(\x-372)/81)*((\x-372)*(\x-372)/81))
    +0.22/(1+((\x-503)*(\x-503)/169)*((\x-503)*(\x-503)/169));
}}
\begin{groupplot}[
  % 【overfull 修正】原为 width=48mm + horizontal sep=8mm：
  % 3×48 + 2×8 = 160 mm，再加 ylabels at=edge left 的左侧轴标签与刻度约 15 mm
  % 共约 175 mm，超出单栏版心 165 mm 约 10 mm ——正是实测的 29.82 pt。
  % 改为 3×45 + 2×6 = 147 mm，含左侧标签约 162 mm，留 3 mm 余量。
  group style={group size=3 by 1, yticklabels at=edge left, ylabels at=edge left,
               xlabels at=edge bottom, horizontal sep=6mm},
  width=45mm, height=38mm,
  xmin=40, xmax=620,
  xtick={100,300,500}, minor xtick={200,400,600},
  xlabel={温度 $T$ / $\degC$},
  scale only axis
]
\nextgroupplot[ymin=46, ymax=108, ytick={60,80,100},
  ylabel={质量 $m$ / \%}]
\addplot[domain=40:620,samples=400,mainline,black]{tgm(x)};
\draw[accentsoft,dashed,thinline] (axis cs:112,46) -- (axis cs:112,108);
\draw[accentsoft,dashed,thinline] (axis cs:258,46) -- (axis cs:258,108);
\draw[accentsoft,dashed,thinline] (axis cs:318,46) -- (axis cs:318,108);
\draw[accentsoft,dashed,thinline] (axis cs:372,46) -- (axis cs:372,108);
\draw[accentsoft,dashed,thinline] (axis cs:503,46) -- (axis cs:503,108);
\node[anchor=south,color=accent] at (axis cs:112,103) {①};
\node[anchor=south,color=accent] at (axis cs:258,103) {②};
\node[anchor=south,color=accent] at (axis cs:345,103) {③};
\node[anchor=south,color=accent] at (axis cs:503,103) {④};

\nextgroupplot[ymin=-0.70, ymax=0.10, ytick={0,-0.4},
  ylabel={$\mathrm{d}m/\mathrm{d}t$ / (\%\,$^{\circ}$C$^{-1}$)}]
\addplot[domain=40:620,samples=400,mainline,black]{dtm(x)};
\draw[accentsoft,dashed,thinline] (axis cs:112,-0.70) -- (axis cs:112,0.10);
\draw[accentsoft,dashed,thinline] (axis cs:258,-0.70) -- (axis cs:258,0.10);
\draw[accentsoft,dashed,thinline] (axis cs:318,-0.70) -- (axis cs:318,0.10);
\draw[accentsoft,dashed,thinline] (axis cs:372,-0.70) -- (axis cs:372,0.10);
\draw[accentsoft,dashed,thinline] (axis cs:503,-0.70) -- (axis cs:503,0.10);
\node[anchor=west,color=accent] at (axis cs:60,-0.40) {峰朝下＝失重};

\nextgroupplot[ymin=-0.75, ymax=0.95, ytick={-0.5,0,0.5},
  ylabel={热流率 $\mathrm{d}Q/\mathrm{d}t$ / (mW\,mg$^{-1}$)}]
\addplot[domain=40:620,samples=400,mainline,black]{dsm(x)};
\draw[accentsoft,dashed,thinline] (axis cs:112,-0.75) -- (axis cs:112,0.95);
\draw[accentsoft,dashed,thinline] (axis cs:258,-0.75) -- (axis cs:258,0.95);
\draw[accentsoft,dashed,thinline] (axis cs:318,-0.75) -- (axis cs:318,0.95);
\draw[accentsoft,dashed,thinline] (axis cs:372,-0.75) -- (axis cs:372,0.95);
\draw[accentsoft,dashed,thinline] (axis cs:503,-0.75) -- (axis cs:503,0.95);
\node[anchor=west,color=accent] at (axis cs:55,0.80) {吸热（$+$）};
\node[anchor=west,color=accent] at (axis cs:55,-0.68) {放热（$-$）};
\node[anchor=north east,color=accent] at (axis cs:612,-0.28) {③ 内叠加放热峰};
\end{groupplot}
\end{tikzpicture}
\caption{同一模型样品上 TG、DTG 与 DSC 曲线的对应关系。读者应看三处：①—④ 的竖虚线把同一
温度区间的三种信号连起来；TG 的四个台阶与 DTG 的四个峰一一对应（峰的起止点即台阶的起止点，
峰位对应台阶的拐点）；而 DSC 上第 ③ 温区由 $\qty{318}{\celsius}$ 附近的宽吸热与
$\qty{372}{\celsius}$ 的放热叠加而成——一个台阶不等于一个反应。纵轴方向约定：TG 向上为
增重、DTG 失重峰朝下；DSC 按热力学惯例向上为正（吸热）、向下为负（放热），DTA 的方向与
此相反。曲线为模型计算：TG 的台阶用有界 S 型函数
$S(T)=\frac{1}{2}\left[1+(T-T_i)/\sqrt{(T-T_i)^2+w_i^2}\right]$ 按
$(\Delta m_i,\,T_i,\,w_i)=(8.0\%,\,110,\,12)$、$(14.0\%,\,255,\,14)$、$(22.0\%,\,355,\,20)$、
$(3.0\%,\,500,\,14)$ 叠加，DTG 为其解析导数，DSC 用 $A/[1+((T-T_i)/w)^4]$ 叠加并加一条
$\qty{0.02}{}\to\qty{0.22}{mW.mg^{-1}}$ 的线性漂移基线。样品与全部数值均为教学设定，不代表
任何实际催化剂。\modelmark\figsrc{信号约定与读图规则据 S1（上册）图 3-2（书内 p.175 /
PDF p.194）与图 3-7（书内 p.179 / PDF p.198）自绘。}}
\label{fig:ch03-tgdsc}
\end{figure}

\begin{dataflow}
（1）先确认纵轴约定：这张图是 TG、DTG 还是 DSC／DTA，吸热朝上还是朝下。（2）用 DTG 划分
台阶边界——DTG 峰的起止点就是 TG 台阶的起止点，峰的数目等于台阶的数目。（3）逐台阶计算
失重百分数，基线取该台阶\textbf{之前}的平台，不要用总失重代替单台阶失重。（4）用 DSC 判断
每个台阶的热效应方向，并检查同一温区内是否叠加了方向相反的事件。（5）把台阶失重换算成
物质的量，与元素分析、化学式或负载量对照，检验归属是否自洽。（6）归属只有在逸出气体被
确认之后才成立——TG 只能告诉你有质量损失，不能告诉你逸出的是什么，须与质谱、红外等联用
或离线分析配合（见 §3.7）。（7）报告时给出升温速率、气氛、样品量与坩埚信息。
\end{dataflow}

\paragraph{一个完整的读图实例}以图~\ref{fig:ch03-tgdsc} 为例。
\pedset{以下样品与全部数值均为教学设定的模型计算值，不是实验数据}。

\begin{obs}
TG 曲线有四个台阶：$60$--$\qty{160}{\celsius}$ 失重 $8.0\%$，$190$--$\qty{300}{\celsius}$
失重 $14.0\%$，$300$--$\qty{430}{\celsius}$ 失重 $22.0\%$，$430$--$\qty{570}{\celsius}$
失重 $3.0\%$，$\qty{570}{\celsius}$ 以上为平台，残余量 $53.0\%$。DTG 上四个峰，峰温分别约
为 $110$、$255$、$355$、$\qty{500}{\celsius}$，其中第三个峰最宽、最高。DSC 上第 ①、②、
④ 台阶各对应一个向上的峰（约 $112$、$258$、$\qty{503}{\celsius}$）；第 ③ 台阶所在温区
信号复杂，$\qty{318}{\celsius}$ 附近是一个宽而低的向上峰，$\qty{372}{\celsius}$ 处叠加
一个明显向下的尖峰。
\end{obs}

把台阶高度换算成物质的量，可以对归属做一次独立的算术检验。设元素分析给出该样品含 Ni
$12.0\%$，则每 $\qty{100}{g}$ 样品中 Ni 为 $12.0/58.69=\qty{0.204}{mol}$；第二台阶失水
$14.0/18.02=\qty{0.777}{mol}$，二者之比 $0.777/0.204=3.8$，接近 $4$，与水合盐或层间水的
化学计量自洽。同一算法给出第一台阶的水量比为 $2.2$，不呈整数比，更符合表面物理吸附水的
特征——这一步不是新的测量，而是用质量数据对归属做了一次一致性检查。

\begin{infer}
按 DSC 的方向约定（向上吸热、向下放热），① ② ④ 三个台阶对应吸热过程，第 ③ 台阶温区内
同时存在吸热与放热两个事件：$\qty{318}{\celsius}$ 附近的宽吸热与盐类分解相符，
$\qty{372}{\celsius}$ 的放热与残留有机物的氧化燃烧相符，而它们在 TG 上合成了一个台阶。
据此判断该样品为负载型金属盐前驱体，其焙烧温区应覆盖第 ③ 台阶的终点。归属仍属推断：
TG 与 DSC 都不能确认逸出气体的种类，须由 TG-MS 或 TG-FTIR 给出 $m/z=18$、$30$、$44$ 的
信号才能定案（见 §3.7）；物相变化须由多晶 X 射线衍射确认（\chapref{ch:04}）。
\end{infer}

\subsection{误判与边界}

\begin{pitfall}
\textbf{（1）升温速率改变峰温。}对动力学控制的过程（脱水、分解），升温速率越高，峰温越向
高温移动、峰越高越窄；对熔化这类热力学控制的过程，峰位移动小，但峰形会因热滞后展宽。
因此不同升温速率测得的峰温不可直接比较，峰温只在升温速率、气氛与装样条件完全相同时才可
作相对比较 \src{S1 §3.2.2，书内 p.177 / PDF p.196}。\textbf{（2）样品量引入传热滞后。}
样品量越大，试样内部温度梯度越大，吸热过程还会使试样实际温度低于程序温度；结果是峰变宽、
峰温偏高、相邻事件合并，定量结果一律偏向「反应更慢、更晚」。用小样品、薄层与良好接触可以
减小但不能消除这一效应，比较不同样品时必须使用同一质量与同一装样方式。\textbf{（3）基线
漂移改变峰面积。}试样与参比的热容不等、传感器不对称、升温本身引起的热流变化，都会使基线
倾斜或弯曲。峰面积是峰与\textbf{内插基线}之间的面积，基线选错会系统性地改变 $\Delta H$；
基线必须由同条件下的空白测量确定，不能在两段平台之间随手连线。\textbf{（4）浮力与气流改变
TG 的表观质量。}升温时炉内气体密度下降、浮力减小，热天平会记录到一个表观增重；气流速率
改变、气氛切换（例如 \ce{N2} 切换为 \ce{O2}）、吊丝与支架的热膨胀同样会造成表观台阶或
漂移。这些效应与样品无关，必须用空坩埚在同程序下测空白并扣除，否则会把基线漂移误读成
氧化增重或吸附。\textbf{（5）坩埚材料与气氛参与反应。}铂坩埚对有机物氧化有催化作用，且会
被卤素、硫、磷等腐蚀；铝坩埚熔点低，高温下与碱性或酸性样品反应；石英坩埚与碱金属氧化物
反应；刚玉坩埚在强还原气氛或金属蒸气下也可能被侵蚀。气氛一侧同样要当心：空气中的 \ce{O2}
会参与氧化，\ce{N2} 在高温下可与某些金属生成氮化物而表现为增重。选择坩埚的原则是「在所测
温区内对样品与气氛都惰性」，并在报告中写明材料。\textbf{（6）台阶数不等于反应数。}$N$ 个
台阶对应 $N$ 个 DTG 峰，但同一个台阶里可以包含方向相反的多个事件（图~\ref{fig:ch03-tgdsc}
的第 ③ 台阶），只凭 TG 判断「有几个反应」会系统性偏少。\textbf{（7）MTDSC 的分量划分不是
材料固有的二分。}可逆与不可逆的归属依赖基础升温速率、调制周期与调制幅度；报告分量曲线时
必须同时给出这三个参数，并与普通 DSC 的总热流对照。\textbf{（8）信号方向不能靠习惯判断。}
DTA 向上为放热、DSC 向上为吸热，两者的约定相反；同一份数据换一台仪器读，若只凭「峰朝上
就是放热」的记忆，会把分解与燃烧的归属整体颠倒。
\end{pitfall}

\textbf{共同的边界。}以上八条中，（1）（2）改变的是温度坐标，（3）（4）改变的是纵坐标，
（5）改变的是样品本身，（6）（7）（8）改变的是读图规则。四类误差的量级不同，但有一个共同
的后果：它们都不会让曲线「看起来出错」。因此热分析的结论必须以\textbf{同一条件下的空白
测量与重复测量}为对照，并以独立方法确认归属，而不是仅凭一条曲线的形状。

\subsection{互补与回查}

\begin{crosscheck}
热分析给出的是「质量—温度」与「热量—温度」两条线，二者互补，但都不能单独定案。
\textbf{（1）质量与热量互相校验。}若某台阶被归属为碳酸盐分解，则由台阶失重算出的碳酸盐
物质的量、与由 DSC 峰面积算出的吸热焓，应彼此自洽；两者矛盾时优先怀疑归属或基线，而不是
取平均值。\textbf{（2）逸出气体必须独立确认。}TG-MS、TG-FTIR 把质量损失与气体分子一一对应，
是台阶归属最直接的证据，详见 §3.7。\textbf{（3）物相与表面状态用别的方法确认。}分解前后的
物相变化由多晶 X 射线衍射给出（\chapref{ch:04}），表面元素与化学态由 X 射线光电子能谱给出
（\chapref{ch:09}）。\textbf{（4）活性与表面性质用程序升温技术确认。}还原性、酸性位与活性位
数量由化学吸附与程序升温还原／脱附给出（\chapref{ch:05}）；比表面与孔结构随焙烧温度的变化
由物理吸附给出（\chapref{ch:01}）。\textbf{（5）动力学参数要与反应动力学实验对照。}由热分析
得到的 $E_{\mathrm{a}}$ 是程序升温、非均相条件下的表观值，与在反应条件下测得的本征动力学
参数不是同一个量（\chapref{ch:10}）；两者数值接近只能作为佐证，不能互相替代。
\end{crosscheck}

\srcnote{本分片来源：S1（上册）§3.1--§3.3，书内 p.169--181 / PDF p.188--200。
式~\eqref{eq:ch03-alpha}--\eqref{eq:ch03-integral} 已逐式对照原页面图像复核（2026-10-13）：
其中式~\eqref{eq:ch03-alpha}、式~\eqref{eq:ch03-dta}、式~\eqref{eq:ch03-dsc}、
式~\eqref{eq:ch03-arrh} 与式~\eqref{eq:ch03-integral} 原书未印，属改编补充（每一式的原书对照
就写在正文相应位置）；式~\eqref{eq:ch03-rate} 的右半与原书
$\mathrm{d}\alpha/\mathrm{d}T=(1/\beta)k(T)f(\alpha)$ 逐字同形；式~\eqref{eq:ch03-mtdsc}
的第二式与原书 $\mathrm{d}H/\mathrm{d}t=C_p\,\mathrm{d}T/\mathrm{d}t+f(T,t)$ 同形而符号不同
（原书记热流率为 $H$，本稿记 $Q$）。逐式判决见 \texttt{sources/formula\_review.md} 的
ch03-a 一节；原书 p.176（PDF p.195）整页无 OCR 文本，该页所载 DTA 曲线特征量定义与配图已按
页面图像补齐。本章其余节次（§3.4 应用、§3.5 联用技术、§3.6 实验技巧、§3.7 结束语）由后续
分片续写。}
